% !TeX root = ../main.tex
\section{Presentación del argumento}
El desarrollo de un ensayo reúne razones que sostienen una tesis y responde a las objeciones relevantes.
Cada párrafo debe cumplir una función reconocible, mantener la continuidad con el anterior y evitar afirmaciones sin respaldo.
Las fuentes se documentan mediante notas al pie y se reúnen al final en la bibliografía.

\subsection{Notas y referencias}
La primera referencia bibliográfica de ejemplo incluye los datos completos de la obra.
Los números de página de este modelo sirven únicamente para demostrar la sintaxis de citación.\autocite[202]{donner2020}
La siguiente llamada muestra la forma abreviada de una obra ya citada.\autocite[293]{donner2020}

\subsubsection{Lectura de las fuentes}
Antes de presentar un ensayo, corresponde verificar directamente la edición consultada, la transcripción de cada cita y su localizador.
La automatización del formato no reemplaza esa revisión.
Una referencia con dos autores ofrece un segundo ejemplo del aparato bibliográfico.\autocite[16]{sendek2007}

\subsubsection{Citas breves}
El siguiente texto se ha redactado para ilustrar las comillas latinas, la puntuación y la llamada de nota: \enquote{Una cita precisa conserva las palabras de su fuente y permite al lector reconocer su procedencia}.\footnote{Texto original de demostración; no es una cita tomada de las obras de ejemplo.}
En el ensayo definitivo, una cita breve se integra en el párrafo y se acompaña de la referencia correspondiente.

\subsection{Citas extensas}
El siguiente bloque contiene más de cuarenta palabras y muestra el tratamiento tipográfico de una cita extensa.
También es texto original de demostración y no reproduce un pasaje de las obras citadas:
\begin{SRLblockquote}
Una exposición teológica responsable permite distinguir entre la afirmación que desea defender, las razones que ofrece para sostenerla y las fuentes que consulta durante su elaboración.
El lector debe poder seguir ese recorrido sin confundir la voz del autor con la de sus interlocutores, ni tomar una ilustración del formato por una evidencia del argumento.\footnote{Bloque original de demostración, incluido exclusivamente para comprobar sangría e interlineado.}
\end{SRLblockquote}

\section{Revisión del documento}
La revisión final comprende el argumento, el uso de las fuentes y la presentación material.
Una obra editada proporciona un tercer ejemplo bibliográfico.\autocite[432]{tuggy1996}
Los encabezados, la numeración y la bibliografía deben revisarse en el PDF generado, especialmente después de cambiar la familia tipográfica.
